\documentclass[11pt]{article}
\usepackage[margin=1in]{geometry}
\usepackage{amsmath,amssymb}

\title{MATH 375: Probability Theory}
\author{Mariana Restrepo}
\date{Fall 2026}

\begin{document}
\maketitle

\begin{center}
\large Lecture 5
\end{center}

\tableofcontents

\section{Lecture 5: Multiplication, Total Probability, and Independence}

\subsection{The multiplication rule}

For events $A$ and $B$, provided $P(A)>0$,
\[
P(A\cap B)=P(A)P(B\mid A).
\]

The rule extends to several events. For example,
\[
\begin{aligned}
P(E_1\cap E_2\cap E_3)
={}&P(E_1)\,
P(E_2\mid E_1)\,
P(E_3\mid E_1\cap E_2).
\end{aligned}
\]
In general, each new factor is conditioned on all the events
already listed before it.

\subsection{Example: separating four aces}

A shuffled deck is divided into four labeled piles of 13 cards.
What is the probability that every pile contains exactly one ace?

Imagine placing the four aces into piles one at a time. The first
ace can go into any pile. For the other three aces to go into
different piles, their choices must differ from the earlier ones:
\[
\begin{aligned}
P(\text{one ace in each pile})
&=1\cdot\frac34\cdot\frac24\cdot\frac14\\
&=\boxed{\frac{3}{32}}.
\end{aligned}
\]

\subsection{The law of total probability}

Suppose $F_1,\ldots,F_n$ are mutually exclusive events whose union
is the whole sample space. They form a \textbf{partition}: exactly
one of them occurs.

Any event $E$ can be split into disjoint pieces:
\[
E=(E\cap F_1)\cup\cdots\cup(E\cap F_n).
\]
Adding the probabilities of those pieces gives
\[
\boxed{
P(E)=\sum_{i=1}^{n}P(F_i)P(E\mid F_i).
}
\]

\subsection{Example: choosing a coin}

A box contains two type A coins and one type B coin. A type A coin
lands heads with probability $\tfrac14$, while a type B coin lands
heads with probability $\tfrac34$. One coin is selected uniformly
and flipped. Given that it lands heads, what is the probability
that it was type A?

Let $A$ mean ``type A was selected,'' $B$ mean ``type B was
selected,'' and $H$ mean ``the flip is heads.'' Then
\[
P(A)=\frac23,\qquad P(B)=\frac13,
\]
and
\[
P(H\mid A)=\frac14,\qquad P(H\mid B)=\frac34.
\]

First, total probability gives
\[
\begin{aligned}
P(H)
&=P(A)P(H\mid A)+P(B)P(H\mid B)\\
&=\frac23\cdot\frac14+\frac13\cdot\frac34\\
&=\frac16+\frac14=\frac5{12}.
\end{aligned}
\]
Now apply Bayes' formula:
\[
\boxed{
P(A\mid H)
=\frac{P(A)P(H\mid A)}{P(H)}
=\frac{1/6}{5/12}
=\frac25.
}
\]

\subsection{Example: at least one girl}

Assume each child is independently equally likely to be a boy or
a girl. For a two-child family, use ordered outcomes:
\[
S=\{BB,BG,GB,GG\}.
\]
Given that \emph{at least one child is a girl}, the remaining
equally likely possibilities are $BG$, $GB$, and $GG$. In only one
of these are both children girls. Therefore,
\[
\boxed{
P(\text{both girls}\mid\text{at least one girl})=\frac13.
}
\]
The wording of the given information determines which outcomes
remain in the sample space.

\subsection{Example: a girl born on Tuesday}

Assume birth sex and weekday are independent and equally likely.
Suppose we are told that \emph{at least one of two children is a
girl born on Tuesday}. What is the probability that both children
are girls?

For one child, the probability of \emph{not} being a girl born on
Tuesday is $13/14$. Thus,
\[
P(\text{at least one girl born Tuesday})
=1-\left(\frac{13}{14}\right)^2
=\frac{27}{196}.
\]

For both children to be girls and at least one to have been born
on Tuesday, consider their two birth weekdays. There are $7^2$
pairs of weekdays, of which $6^2$ have no Tuesday:
\[
P(\text{both girls and at least one Tuesday girl})
=\frac14\left(1-\frac{6^2}{7^2}\right)
=\frac{13}{196}.
\]
Therefore,
\[
\boxed{
P(\text{both girls}\mid\text{at least one Tuesday girl})
=\frac{13/196}{27/196}
=\frac{13}{27}.
}
\]

\subsection{Independent events}

Events $A$ and $B$ are \textbf{independent} if
\[
\boxed{P(A\cap B)=P(A)P(B).}
\]
If $P(A)>0$, this is equivalent to
\[
P(B\mid A)=P(B).
\]
Knowing whether $A$ occurred does not change the probability of $B$.

\paragraph{Example: two dice.}
Let $F$ mean ``the first die shows 4,'' so $P(F)=1/6$.

Let $E_5$ mean ``the sum is 5.'' There are five ways to get a sum
of 5, and only $(4,1)$ also belongs to $F$. Thus,
\[
P(E_5)=\frac5{36},\qquad
P(E_5\cap F)=\frac1{36},
\]
but
\[
P(E_5)P(F)=\frac5{216}\ne\frac1{36}.
\]
So $E_5$ and $F$ are \textbf{dependent}.

Now let $E_7$ mean ``the sum is 7.'' There are six ways to get
a sum of 7, and only $(4,3)$ also belongs to $F$:
\[
P(E_7)=\frac6{36}=\frac16,\qquad
P(E_7\cap F)=\frac1{36}.
\]
This time,
\[
P(E_7)P(F)=\frac16\cdot\frac16=\frac1{36},
\]
so $E_7$ and $F$ are \textbf{independent}.

\end{document}